\documentclass{article}
\usepackage[utf8]{inputenc}
\usepackage{a4wide}

\begin{document}

\section*{แต้มพลังนักวิจัย}

กำหนดให้\texttt{Citations}เป็นอาร์เรย์ของจำนวนเต็มโดยที่\texttt{citations{[}i{]}}หมายถึงจำนวนการ\textbf{อ้างถึง}บทความวิจัยฉบับที่\texttt{i}ของศาสตราจารย์ต๋อย

ให้\textbf{คำนวณหาพลัง}\texttt{M}\textbf{ของศาสตราจารย์ต๋อย}ซึ่งค่า\texttt{M}หมายถึงค่าสูงสุดของ\texttt{n}โดยที่นักวิจัยคนหนึ่งมีงานวิจัยอย่างน้อย\texttt{n}เรื่องที่ทุกเรื่องในนั้นถูกอ้างถึง\textbf{อย่างน้อย}\texttt{n}ครั้ง\\



\hfill\break



{\textbf{แหล่งข้อมูลที่น่าสนใจ:}}{\url{https://en.wikipedia.org/wiki/H-index}}



\hfill\break



\textbf{ตัวอย่าง}



\begin{quote}

\textbf{Input:}



5



3 0 6 1 5



\textbf{Output:}3



\texttt{3\ 0\ 6\ 1\ 5}หมายถึง{ศาสตราจารย์ต๋อย}{มีบทความทั้งหมด 5 เรื่อง

และได้รับการอ้างอิง\textbf{3, 0, 6, 1, 5}เรื่อง ตามลำดับ}



{{เนื่องจาก{ศาสตราจารย์ต๋อย}มีบทความ 3 เรื่องที่มีการอ้างอิงอย่างน้อย 3 เรื่องต่อเรื่อง

และอีก 2 เรื่องที่เหลือมีเรื่องอ้างอิงไม่เกิน 3 เรื่องต่อเรื่อง

ดังนั้นค่าพลังของ{ศาสตราจารย์ต๋อย}จึงเป็น 3}\\

}

\end{quote}



\subsection*{Input}
บรรทัดแรก จำนวนงาน{วิจัยของ}{ศาสตราจารย์ต๋อย}



บรรทัดที่ 2 งานวิจัยของ{ศาสตราจารย์ต๋อย คั่นด้วยช่องว่าง 1 ช่อง}



\subsection*{Output}
{พลัง}\texttt{M}{ของศาสตราจารย์ต๋อย}\\



\end{document}
